\section{A fixed cubic Selmer invariant and an integral CM coefficient realization}
\label{sec:cm-twisted-selmer}

This section advances two parts of the remaining supply problem without
repeating the Shapiro and Fourier-lattice comparisons of ALP\_MACHINE,
AM-H06--AM-H08.  The first result computes the exact residual strict
Selmer group for the fixed cubic character at the coefficient prime
$3$.  The second realizes that character by an integral Hom lattice
between two elliptic curves over the fixed CM base.  Neither result
assumes A, and neither asserts the arithmetic-window bound (H04.4).

\subsection{The exact residual strict Selmer group}

Put
\[
 K=\mathbb Q(\omega),\quad \omega^2+\omega+1=0,\quad
 E=K(\alpha),\quad \alpha^3=269,\quad
 \chi(g)=g(\alpha)/\alpha\in\mu_3.
\]
Let $\mathcal O=\mathbb Z_3[\omega]$, $T=\mathbb Q_3(\omega)$,
$\lambda=\omega-1$, and write $\varepsilon_3$ for the $3$-adic
cyclotomic character.  The set $S$ consists of the unique primes of
$K$ above $3$ and $269$.  For any integer $e$ let
\[
 M_e=(T/\mathcal O)(\chi\varepsilon_3^{-e})
\]
and define the strict group using exactly the local conditions of
AM-H06:
\[
 \mathcal S_e=\ker\left(
 H^1(G_{K,S},M_e)\longrightarrow
 H^1(G_{K_3},M_e)\oplus H^1(I_{269},M_e)\right).
\]
Here $K_3=\mathbb Q_3(\omega)$, and the second term is inertia at the
prime above $269$.  All cohomology is continuous.  In the notation of
AM-H08, whenever the length is finite,
$t_{3,\chi}(e)=\operatorname{length}_{\mathcal O}\mathcal S_e$.

\begin{theorem}[AC-G22: a zero group at weight zero and one residual generator otherwise]
\label{AC-G22}
The group $\mathcal S_0$ is zero.  For every nonzero integer $e$,
\[
 \mathcal S_e[\lambda]\simeq\mathbb F_3.
\]
Its nonzero residual line is the image of the Kummer class of
$\omega$ under
$H^1(G_{K,S},\mathbb F_3)\to H^1(G_{K,S},M_e)$, after identifying
$\mathbb F_3$ with $\mu_3$.

More precisely, for $e\ne0$ the $\mathcal O$-module $\mathcal S_e$
is either $\mathcal O/\lambda^{m_e}\mathcal O$ for an integer
$m_e\geq1$, or $T/\mathcal O$.  For every positive integer $e$,
the first alternative holds unconditionally.  Thus throughout the
positive-weight range of AM-H04 the entire individual twisted
length is one finite exponent $m_e$.  This theorem does not bound
that exponent.  For negative $e$ the divisible alternative is not
excluded by the residual calculation.
\end{theorem}
\begin{proof}
\proofstep{1. The residual global and local Kummer spaces.}
The ring $\mathbb Z[\omega]$ is norm-Euclidean and has units $\mu_6$.
Consequently the $S$-unit group is
\[
 \mathcal O_{K,S}^{\times}
   =\mu_6\times\langle\lambda\rangle\times\langle269\rangle.
\]
Kummer theory, together with the trivial class group, gives
\[
 H^1(G_{K,S},\mathbb F_3)
   =\mathbb F_3[\omega]\oplus
     \mathbb F_3[\lambda]\oplus
     \mathbb F_3[269].                                  \tag{G22.1}
\]
For clarity, an unramified-outside-$S$ Kummer class has valuations
divisible by $3$ at all primes outside $S$; principality removes those
valuations by a global cube.  Conversely an $S$-unit is unramified at
every prime outside $S$, since its residue characteristic is not $3$
and the Kummer polynomial has unit derivative at its roots.  This
proves the stated global description.  It is the same elementary
Kummer calculation used for dimensions in AM-H07, now with its
generators retained.

The integer $269$ is a cube already in $\mathbb Q_3$.  Indeed
\[
 f(t)=t-3t^2+3t^3-30
\]
has $f(0)\equiv0\pmod3$ and $f'(t)\equiv1\pmod3$.
Hensel lifting gives $t\in\mathbb Z_3$ with $f(t)=0$, and
$(-1+3t)^3=-1+9\cdot30=269$.
Thus $\chi$ is trivial on $G_{K_3}$ and $[269]$ restricts to zero
locally.  The two local classes $[\lambda]$ and $[\omega]$ are
independent: the valuation of $\lambda$ is one, whereas $\omega$ is
a unit and is not a cube in $K_3$.  The latter assertion follows from
$[\mathbb Q_3(\zeta_9):\mathbb Q_3(\zeta_3)]=3$; adjoining a cube
root of $\omega$ gives that extension.  Therefore restriction on
(G22.1) has kernel exactly $\mathbb F_3[269]$ and image the
two-dimensional span of $[\lambda],[\omega]$.

\proofstep{2. The ramified inertia condition is automatic for this module.}
At $269$, $E/K$ is a tamely totally ramified cubic extension, by the
Eisenstein polynomial $X^3-269$.  The cyclotomic character
$\varepsilon_3$ is unramified there.  The maximal pro-$3$ quotient
of tame inertia is $\mathbb Z_3$, with a generator acting on $M_e$
by $\omega$.  Multiplication by $\omega-1=\lambda$ is surjective
on $T/\mathcal O$.  Every continuous one-cocycle on this procyclic
group is therefore a coboundary: choose a preimage under $\lambda$
of its value on a generator, subtract that coboundary, and use the
density of the generator's integer powers.  The prime-to-$3$ inertia
factors have vanishing positive-degree cohomology on these
$3$-primary coefficients, by averaging on their finite quotients.
It follows that
\[
 H^1(I_{269},M_e)=0.                                    \tag{G22.2}
\]
In particular the actual ramified coefficient action has been used;
it has not been replaced by an unramified coefficient system.

\proofstep{3. The global connecting line.}
All actions on $M_e[\lambda]$ are trivial.  The exact sequence
\[
 0\longrightarrow\mathbb F_3\longrightarrow M_e
     \xrightarrow{\lambda}M_e\longrightarrow0
\]
gives for any group $H$ in question
\[
 H^1(H,M_e)[\lambda]
  \simeq H^1(H,\mathbb F_3)/D_H,
 \qquad
 D_H=\operatorname{im}\left(
 M_e^H/\lambda M_e^H\xrightarrow{\delta_H}
 H^1(H,\mathbb F_3)\right).                            \tag{G22.3}
\]
Globally inertia at $269$ acts by $\omega$, while every group
element acts as one modulo $\lambda$.  Hence
$M_e^{G_{K,S}}=M_e[\lambda]$.
Use $1/\lambda$ as its generator and $1/\lambda^2$ as a lift through
multiplication by $\lambda$.  With
$\rho_e=\chi\varepsilon_3^{-e}$, the connecting cocycle becomes
\[
 g\longmapsto\frac{\rho_e(g)-1}{\lambda}\pmod\lambda.
\]
Since $\varepsilon_3(g)\equiv1\pmod3$ on $G_K$ and
$3\mathcal O=\lambda^2\mathcal O$, the cyclotomic part contributes
zero to this expression.  If $\chi(g)=\omega^a$, the value is
$a\pmod3$.  It is exactly the additive character cutting out
$K(\sqrt[3]{269})$, so
\[
 D_{G_{K,S}}=\mathbb F_3[269]                           \tag{G22.4}
\]
for every integer $e$.

\proofstep{4. The local connecting line at three.}
For $e=0$ the local action is trivial and
$M_0^{G_{K_3}}=T/\mathcal O$ is divisible; hence
$D_{G_{K_3}}=0$.
For $e\ne0$ the local cyclotomic image is
$1+3\mathbb Z_3$.  If $m=v_3(e)$, the elementary lifting-of-exponent
calculation gives
\[
 \min_{g\in G_{K_3}}v_\lambda
       (\varepsilon_3(g)^{-e}-1)=2+2m=:d_e.
\]
Indeed raising a principal unit to the third power raises its
$3$-adic distance valuation from one by exactly one at every step;
raising to an exponent prime to $3$, or taking an inverse, preserves
that valuation.  The element $4\in1+3\mathbb Z_3$ realizes the
minimum.  Therefore
$M_e^{G_{K_3}}=M_e[\lambda^{d_e}]$.
Lifting its generator $1/\lambda^{d_e}$ to
$1/\lambda^{d_e+1}$ shows that the local connecting line is generated
by
\[
 g\longmapsto
  \frac{\varepsilon_3(g)^{-e}-1}{\lambda^{d_e}}\pmod\lambda.
\]
This is a nonzero scalar multiple of
$(\varepsilon_3(g)-1)/3\pmod3$.
For example, write $e=3^m e_0$ with $3\nmid e_0$ and use
$\lambda^2=-3\omega$ in the binomial expansion.  The character has
kernel $\varepsilon_3\equiv1\pmod9$, and therefore cuts out
$K_3(\zeta_9)/K_3$.  Its Kummer line is $\mathbb F_3[\omega]$.
We have proved
\[
 D_{G_{K_3}}=
 \begin{cases}0,&e=0,\\
 \mathbb F_3[\omega],&e\ne0.
 \end{cases}                                           \tag{G22.5}
\]
The global invariant generator lies inside
$\lambda M_e^{G_{K_3}}$ when $e\ne0$, because $d_e\geq2$;
thus its local boundary is zero, as also follows from the vanishing
of the local class $[269]$.  This checks compatibility of the two
connecting maps.

\proofstep{5. The strict kernel and its module structure.}
Naturality in (G22.3), together with (G22.2), identifies
\[
 \mathcal S_e[\lambda]
 =\frac{\operatorname{res}_3^{-1}(D_{G_{K_3}})}
        {D_{G_{K,S}}}.                                  \tag{G22.6}
\]
If $e=0$, the numerator and denominator are both
$\mathbb F_3[269]$.  If $e\ne0$, the numerator is
$\mathbb F_3[269]\oplus\mathbb F_3[\omega]$ and the denominator
is its first summand.  This proves both residual assertions.

Every continuous cocycle from the compact group $G_{K,S}$ to the
discrete torsion module $M_e$ has finite image.  Thus every
cohomology class, and every element of $\mathcal S_e$, is annihilated
by some power of $\lambda$.  A nonzero such element generates a
finite cyclic module whose socle is nonzero.  Consequently a zero
socle implies $\mathcal S_0=0$.

Finally a $\lambda$-primary torsion module with one-dimensional socle
has a particularly simple form.  Every finitely generated submodule
is finite; the structure theorem over the DVR $\mathcal O$ and its
one-dimensional socle make it cyclic.  Any two cyclic submodules
are comparable, since the module they generate is cyclic and its
submodules are powers of $\lambda$.  If their lengths are bounded,
the entire module is one finite cyclic module.  If unbounded, their
union is the unique divisible rank-one torsion module
$T/\mathcal O$: choose compatible generators of successively larger
cyclic submodules.  This proves the stated dichotomy.

\proofstep{6. Finiteness at every positive weight.}
Here the additional input is a proved theorem for arbitrary number
fields, not a consequence of the one-dimensional socle.  In
C.~Soul\'e, \emph{\'El\'ements cyclotomiques en K-th\'eorie},
Ast\'erisque 147--148 (1987), \S6.1, printed pp.~246--247,
$H^2(G_{F,S_3},\mathbb Z_3(n))$ is finite for every number field
$F$ and $n\geq2$, and the adjacent coefficient exact sequence gives
$H^2(G_{F,S_3},\mathbb Q_3/\mathbb Z_3(n))=0$.
Here $S_3$ consists precisely of the primes above $3$.
The definition and criterion in \S6.2, printed p.~247, identify
\[
 R_i(F)=\ker\left(H^1(G_{F,S_3},\mathbb Q_3/\mathbb Z_3(i))
       \longrightarrow\prod_{w\mid3}
              H^1(F_w,\mathbb Q_3/\mathbb Z_3(i))\right)
\]
and assert that, for $i\ne0$, this group is finite if and only if
$H^2(G_{F,S_3},\mathbb Q_3/\mathbb Z_3(1-i))=0$.
Taking $i=-e<0$ makes $1-i=e+1\geq2$.
Thus $R_{-e}(E)$ is finite.  This use of the source does not impose
abelian or CM hypotheses on $E$.

For completeness, finiteness passes to our character without
assuming a semisimple integral character decomposition.  Restrict
classes from $K$ to $E$.  On $G_E$ the coefficient $M_e$ is
$(T/\mathcal O)(-e)$, which as a $\mathbb Z_3$-module is a direct
sum of two copies of $\mathbb Q_3/\mathbb Z_3(-e)$.
The restriction of a strict class is zero at the primes above $3$
and on inertia at the primes above $269$.  Since the coefficient
is unramified outside $3$ after restriction, the latter condition
means the class descends from $G_{E,S_3}$: by inflation--restriction
its restriction to the normal subgroup generated by those inertia
groups is zero.  Its image therefore belongs to the finite group
$R_{-e}(E)\otimes_{\mathbb Z_3}\mathcal O$.

The kernel of restriction injects into
$H^1(\operatorname{Gal}(E/K),M_e^{G_{E,S}})$ by
inflation--restriction.  The invariant coefficient group is finite:
the cyclotomic image on $G_E$ is an open subgroup of
$\mathbb Z_3^{\times}$, and for $e>0$ it contains an element
$u$ with $u^e\ne1$, so these invariants are killed by the nonzero
scalar $u^{-e}-1$.  Cohomology of the finite group
$\operatorname{Gal}(E/K)$ with a finite coefficient group is finite.
Both the image and kernel of restriction are consequently finite,
and so is $\mathcal S_e$.  The already proved one-dimensional
socle identifies it with $\mathcal O/\lambda^{m_e}$ for some
$m_e\geq1$.
\end{proof}

\begin{corollary}[Untwisted base term in AM-H08]
\label{AC-G22-base}
For every integer $e$, the strict group
\[
 \ker\left(H^1(G_{K,S_3},\mathbb Q_3/\mathbb Z_3(-e))
       \longrightarrow H^1(K_3,\mathbb Q_3/\mathbb Z_3(-e))\right)
\]
is zero.  In particular the base length $a_K(e)$ in AM-H08 is zero.
\end{corollary}
\begin{proof}
Now $S_3$ has only one prime, so its global residual Kummer space
is the span of $[\omega],[\lambda]$, which injects into the local
Kummer space by step 1 above.  Both global and local cyclotomic
images are $1+3\mathbb Z_3$; globally this follows from the
cyclotomic tower $\mathbb Q(\zeta_{3^n})/K$, and locally from
the same tower of completions.  If $e\ne0$, the connecting lines
for multiplication by $3$ on $\mathbb Q_3/\mathbb Z_3(-e)$ are
both the Kummer line $[\omega]$, by the calculation in step 4
with $v_3$ in place of $v_\lambda$.  If $e=0$, both lines are
zero.  In either case the induced map on the residual quotients
is injective.  The strict group therefore has zero $3$-torsion,
and a $3$-primary torsion group with that property is zero.
\end{proof}

\begin{remark}[The remaining quantitative invariant]
For positive $e$, the exponent $m_e$ is equivalently the
largest integer $n$ for which its canonical nonzero socle class lifts
through $\lambda^{n-1}$ inside the \emph{strict} Selmer group.
Thus the missing individual estimate is now an estimate for the
depth of one specific class.  A bound on residual dimension does not
bound that depth: the modules $\mathcal O/\lambda^n$ have the same
one-dimensional socle for every $n\geq1$, and $T/\mathcal O$ has that
socle as well.  These are counterexamples to the formal inference
from socle dimension to length, not assertions that every such
module occurs as one of the arithmetic groups above.
The vanishing at $e=0$ is an actual bound, but the source window
uses positive weights, so it is not a closure of (H04.4).
\end{remark}

\subsection{The character in an integral Hom of two CM elliptic curves}

\begin{proposition}[AC-G23: integral CM cubic twist and its exact elementary costs]
\label{AC-G23}
For a rational prime $p$ put $\alpha^3=p$,
$E_p=K(\alpha)$, $\chi_p(g)=g(\alpha)/\alpha$, and consider the
elliptic curves over the fixed field $K$
\[
 C_0:y^2=x^3+1,\qquad C_p:y^2=x^3+p^2.
\]
There is an origin-preserving isomorphism over $E_p$
\[
 \varphi:C_0\longrightarrow C_p,
 \qquad (x,y)\longmapsto(\alpha^2x,py).                  \tag{G23.1}
\]
Both curves have full CM by $\mathcal O_K=\mathbb Z[\omega]$,
defined over $K$.  For every rational coefficient prime $\ell$,
put $R_\ell=\mathcal O_K\otimes\mathbb Z_\ell$.  The integral
covariant Tate modules are free of rank one over $R_\ell$, and
\[
 \operatorname{Hom}_{R_\ell}
       (T_\ell C_0,T_\ell C_p)
       \simeq R_\ell(\chi_p^2),\qquad
 \operatorname{Hom}_{R_\ell}
       (T_\ell C_p,T_\ell C_0)
       \simeq R_\ell(\chi_p).                            \tag{G23.2}
\]
These are integral identities, including $\ell=3$; no division by
$3$ or integral Fourier projector is used.  In particular the module
$\mathcal O(\chi)$ in AC-G22 is an actual integral CM-linear Hom
lattice over $K$, and its Tate twists are obtained by the
cyclotomic determinant of an elliptic Tate module.

The rational two-dimensional finite-monodromy representation
underlying this Hom is the augmentation representation of
$\operatorname{Gal}(E_p/K)=C_3$.  It is an Artin submotive of the
appropriate tensor of the two elliptic $h^1$ motives; its two
character lines appear after extending coefficients to contain
$\mu_3$.  Its finite-monodromy weight is zero.  This assertion does
not say that $E_p$ is a CM field.

The following quantitative data are exact in the displayed
coordinates.  Let $\eta_i=dx/(2y)$ on $C_i$, and let $h$ denote
absolute logarithmic projective height.  Then
\[
 \varphi^*\eta_p=\alpha^{-1}\eta_0,
 \quad h(\alpha)=h(\alpha^{-1})=\frac13\log p,
 \quad h([\alpha^2:p:1])=h([\alpha^{-2}:p^{-1}:1])=\log p.
                                                               \tag{G23.3}
\]
On a $d$-fold tensor in the corresponding monomial coordinate bases,
the coordinate-matrix height is at most $d\log p$, and the
differential factor has height $(d/3)\log p$.  The stable Faltings
heights of $C_0$ and $C_p$ are equal.

For $p=269$ the isomorphism (G23.1) is defined and integral over
$K_3$ with unit scaling factors.  At the prime $\mathfrak p=(269)$
of $K$, the displayed model of $C_p$ is minimal with
$v_{\mathfrak p}(\Delta)=4$ and its two-dimensional Tate
representation has conductor exponent $2$.  The rank-one
CM-linear Hom character in (G23.2) has conductor exponent $1$ at
$\mathfrak p$ and is unramified at all other finite primes, as in
AC-K04.
\end{proposition}
\begin{proof}
\proofstep{1. Isomorphism, endomorphisms and Galois descent.}
Substitution in the equations gives
$(py)^2=(\alpha^2x)^3+p^2$, since $\alpha^6=p^2$.
The inverse scales $x$ by $\alpha^{-2}$ and $y$ by $p^{-1}$;
the point at infinity is fixed.  For $a\in\mu_3$ let
$[a](x,y)=(ax,y)$.  These automorphisms are defined over $K$ and
commute with $\varphi$.  In characteristic zero the endomorphism
algebra of an elliptic curve containing $\omega$ is $K$: in a
complex uniformization it consists of scalars preserving a
rank-two lattice and hence has degree at most two over $\mathbb Q$.
The endomorphism order already contains the maximal order
$\mathbb Z[\omega]$, so it equals that order.  All its elements,
being integral combinations of $1$ and $[\omega]$, are defined over
$K$.

Composition with $\varphi$ identifies
$\operatorname{Hom}_{\overline K}(C_0,C_p)$ with
$\operatorname{End}_{\overline K}(C_0)=\mathcal O_K$.
For $g\in G_K$,
\[
 g(\varphi)=[\chi_p(g)^2]\circ\varphi.                 \tag{G23.4}
\]
Thus this rank-one $\mathcal O_K$-module has character
$\chi_p^2$, and its dual has character $\chi_p$.

\proofstep{2. The integral coefficient lattice.}
One can check the Tate-module freeness directly from complex
uniformization.  An elliptic curve with endomorphism ring
$\mathcal O_K$ is analytically $\mathbb C/\Lambda$ with
$\Lambda$ a rank-one projective $\mathcal O_K$-module; after
$\ell$-adic completion it is free over each factor of $R_\ell$.
Its $\ell^n$ torsion is $\ell^{-n}\Lambda/\Lambda$,
so passing to the inverse limit gives a free rank-one $R_\ell$
Tate module.  The primary reference for this integral statement,
including coefficient primes ramified in the CM order, is
H.~W. Lenstra, \emph{Complex Multiplication Structure of Elliptic
Curves}, J. Number Theory 56 (1996), Theorem 3 and its Tate-module
reformulation, printed p.~229; the statement was checked in the
author's supplied-online PDF.

Because $\varphi$ has degree one, its induced map on Tate modules is
an integral $R_\ell$-linear isomorphism.  It generates the rank-one
module of all $R_\ell$-linear maps.  Formula (G23.4) therefore proves
(G23.2) at the level of integral lattices.  At $\ell=3$,
$R_3=\mathcal O$, so taking the dual Hom and then the divisible
quotient gives exactly $(T/\mathcal O)\chi$.
The Weil pairing identifies
$\det_{\mathbb Z_\ell}T_\ell C_0$ with
$\mathbb Z_\ell(1)$; tensoring with its $(-e)$-th power, with
duals for negative powers, supplies the Tate factor in $M_e$.

This integral Hom is not obtained by applying $1/3$ times an
averaging operator to an arbitrary lattice.  It is the directly
defined kernel of the CM commutation equations inside the full Hom.
A rational equivariant projector onto it exists, but its integral
splitting at $3$ is a separate issue and is not needed for (G23.2).

\proofstep{3. The Artin-motive identification and variance.}
Over $E_p$, the two-dimensional rational Hom space is generated by
$\varphi$ and $[\omega]\varphi$.  Galois multiplication by
$\omega^2$ has trace $-1$ on nonidentity elements of $C_3$ and trace
$2$ on the identity.  These are the characters of the augmentation
representation, giving the asserted Artin representation.

The motive can also be identified by algebraic correspondences.
The three conjugate isomorphisms
$\varphi,[\omega]\varphi,[\omega^2]\varphi$ give a map from
$h^0(\operatorname{Spec}E_p)$ to the tensor Hom of the two $h^1$
motives.  On $h^1$ their sum is zero because
$1+[\omega]+[\omega]^2=0$ as elliptic endomorphisms.
Thus the map factors through the augmentation summand.  Under the
Rosati trace pairing their Gram matrix has diagonal entries $2$ and
off-diagonal entries $-1$: it is
$\operatorname{Tr}_{K/\mathbb Q}(a\overline b)$ for
$a,b\in\{1,\omega,\omega^2\}$.  On the augmentation plane this
matrix is $3$ times the identity.  Dividing the transpose map by $3$
therefore gives an algebraic splitting with rational coefficients.  Hence the map
is a split injection of rational motives, not just an equality of
Frobenius traces.

For covariant Tate modules the ambient tensor is
$(T_\ell C_0)^\vee\otimes T_\ell C_p$.
With contravariant etale $H^1$, the maps induced by the same
isomorphisms lie instead in
$H^1(C_p)^\vee\otimes H^1(C_0)$.
Interchanging the curves conjugates the individual character; the
rational two-dimensional augmentation constituent is unchanged.
This fixes the variance and does not identify a finite character
with one of the weight-one elliptic characters themselves.

\proofstep{4. Heights, differentials and the fixed local costs.}
Pulling back the displayed differential gives
$\alpha^2/p=\alpha^{-1}$, proving its formula.  The number
$\alpha$ is an algebraic integer and all its complex conjugates
have modulus $p^{1/3}$; the normalized product-formula definition of
height gives $h(\alpha)=\log p/3$ and
$h(\alpha^{-1})=h(\alpha)$.
The projective matrix coefficients $[\alpha^2:p:1]$ are integral
and have archimedean maximum $p$, so their height is $\log p$.
The inverse coefficient triple is projectively
$[p\alpha^{-2}:1:p]=[\alpha:1:p]$, with the same height.
In the natural tensor coordinates every coefficient is a product
of $d$ such diagonal coefficients; their archimedean maximum is at
most $p^d$ and the finite maxima contribute no additional term after
the same integral normalization.  This proves the tensor bound;
the differential ratio is $\alpha^{-d}$ and has the asserted
height.  The stable Faltings height is invariant under isomorphism
over a finite extension.  Since (G23.1) is such an isomorphism,
the two stable heights are exactly equal.

For $269$, the local cube constructed in AC-G22 lies in
$\mathbb Q_3^\times$ and is a unit.  Both scaling factors and their
inverses are $3$-adic units, proving the local integral isomorphism.
At $\mathfrak p=(269)$, the discriminant is
$-432p^4$, so its valuation is $4$.  It is minimal: a
change to a smaller integral discriminant would decrease this
valuation by a positive multiple of $12$, which is impossible
for the nonnegative valuation of the discriminant of an integral
model.  The curve $C_0$ has good reduction at this prime.  After the
tame cubic extension given by $\alpha$, $C_p$ is isomorphic to it.
On its two-dimensional Tate representation at any auxiliary
coefficient prime $\ell\ne269$, inertia therefore acts
through the two nontrivial eigenvalues $\omega,\omega^2$.
There are no inertia invariants and no Swan term, so the Artin
conductor exponent is $2$.  The rank-one Hom character has just
one nontrivial tame eigenvalue, hence exponent $1$.
The conductor calculation AC-K04, or directly the same Kummer local
calculation at every other prime including the local cube at $3$,
gives conductor exactly $(269)$ for that character.
\end{proof}

\begin{remark}[What this construction changes, and what it does not]
The non-CM nature of the splitting field excludes a literal CM field
containing $E_p$; it does not exclude this CM coefficient
construction over $K$.  The curves have fixed dimension, fixed
geometric CM type and equal stable geometric heights.  Their
nontrivial arithmetic input is the cubic descent datum and its
finite conductor, with the explicit coordinate and differential
costs in (G23.3).  These costs grow only linearly in tensor degree
when measured as logarithmic heights for fixed $p$.

This is not yet a substitution in the archived theta or
$\mathrm{VVle}$ budget.  A change of coordinates in a specified
integral theta lattice, a projector on its coefficient system,
and a norm bound on differentiated evaluations have additional
normalizations.  In particular an estimate linear in logarithmic
height becomes exponential after exponentiation.  Neither stable
Faltings-height equality nor the unitary action of $\mu_3$ makes
those costs vanish.  To use (G23.2) one must show that the relevant
theorem permits this rank-one CM-linear Hom coefficient, preserves
the strict local conditions of AC-G22, and controls its tensor-degree
cost uniformly along the actual growing auxiliary hosts.  That
coefficient-version theorem and the bound on $m_e$ remain explicit
gates.
\end{remark}
